\documentclass[12pt,a4paper]{article}

\usepackage{fontawesome5}
\providecommand{\faPencilAlt}{\faIcon{pencil-alt}}

\usepackage{styles/couleurs}
\usepackage{styles/mise_en_page}
\usepackage{styles/boites}
\usepackage{styles/maths}
\usepackage{styles/figures}
\usepackage{styles/evaluations}

\usepackage{enumitem}
\usepackage{array}
\usepackage{tabularx}
\usepackage{amsmath}
\usepackage{amssymb}
\usepackage{tikz}

\usetikzlibrary{calc,patterns}

\setlist[itemize]{leftmargin=1.4em,itemsep=0.15em,topsep=0.2em}
\setlist[enumerate]{leftmargin=*,itemsep=0.25em,topsep=0.25em}

\begin{document}

\enteteinterro[Calculatrice autorisée]{4e}{Médiane et diagramme circulaire}{18 min}{12}

\exo{4}[Calculer]
Détermine la médiane de chaque série. \pts{4}

\[
A:\quad12\quad7\quad10\quad5\quad9
\]

\[
B:\quad4\quad11\quad8\quad6\quad15\quad9
\]

\lignes{5}

\exo{3}[Communiquer]
La médiane d'une série de durées est \(18\) minutes.
Interprète cette information avec une phrase comportant
« au moins la moitié ». \pts{3}

\lignes{3}

\exo{3}[Calculer]
Dans une classe de \(30\) élèves, \(12\) viennent en bus.

\begin{enumerate}
  \item Calcule la fréquence en pourcentage. \pts{1}
  \item Calcule l'angle du secteur correspondant
        dans un diagramme circulaire. \pts{2}
\end{enumerate}

\lignes{4}

\exo{2}[Raisonner]
Un diagramme circulaire comporte des secteurs de :

\[
90^\circ,\qquad72^\circ,\qquad126^\circ,\qquad62^\circ.
\]

Explique pourquoi il ne peut pas représenter une répartition complète. \pts{2}

\lignes{3}

\end{document}
